

In this section, we show that the functor $\GS_{\zeta}$ defined in \cref{subsec the functor} is well-defined, i.e., the web relations in \cref{fig:webrel} hold in $\dmSBSBim$ after applying $\GS_{\zeta}$.

\subsection{Tag Switching Relations}\label{subsec- tag switching relation}
For any valid region labels, we have the tag switching relation in $\cwebs$ given below (up to mirror images and arrow reversals):
\begin{align}
    {
        \phantom{ \tikz[x=1mm, y=1mm, baseline=-0.5ex]{\draw (0,-10) -- (0,10)}} %for vertical spacing
        \tikz[x=1mm,
  y=1mm,
  baseline=-0.5ex] {
          \draw[->] (0,0) -- (0,6) node [right] {$n-k$};
          \draw[->] (0,0) -- (0,-6) node [right] {$k$};
          \draw (0,0) -- (2,0);
        }
        = (-1)^{k(n - k)} \tikz[x=1mm, y=1mm, baseline=-0.5ex]{
          \draw[->] (0,0) -- (0,6) node [right] {$n-k$};
          \draw[->] (0,0) -- (0,-6) node [right] {$k$};
          \draw (0,0) -- (-2,0);
        }
      }. \label{diag-tag switching}\end{align}
We have that
\[ \isvg{0.15}{tag_right_d} = \isvg{0.15}{tag_left_d_rotated}   \]
The LHS of \cref{diag-tag switching} is sent via $\GS_{\zeta}$ to
$$(-1)^{k(n-k)} \nu_{n-k}\; \isvg{0.15}{image_of_tag_rotated}= (-1)^{k(n-k)}\nu_{n-k}\; \isvg{0.15}{image_of_tag_d}.$$
where the last equality is from the isotopy relations.
The RHS of \cref{diag-tag switching} is sent via $\GS_{\zeta}$ to 
$$(-1)^{k(n-k)} \nu_{k}\; \isvg{0.15}{image_of_tag_d}.$$
Hence, we see that \cref{diag-tag switching} holds after applying $\GS_{\zeta}$ if and only if $\nu_k=\nu_{n-k}$.


Similarly, for the relation \cref{diag-tag switching} with reverse arrows, we have the condition that $\lambda_k=\lambda_{n-k}$.
These conditions are exactly those in \cref{cond-tag switching}.
There are no non-trivial mirror images to consider in this case.

\subsection{Tag Cancellation}\label{subsec-tag cancellation}
For any valid region labels, we have the tag cancellation relation in $\cwebs$ given below (up to mirror images and arrow reversals):
\begin{align}
    {
        \tikz[x=1mm,
         y=1mm,
        baseline=-0.5ex] {
          \draw[->-] (0,-6) -- (0,-2);
          \draw (0,-2) -- (2,-2);
          \draw[->-] (0,2) -- (0,-2);
          \draw (0,2) -- (2,2);
          \draw[->] (0,2) -- (0,6) node [right] {$k$};
        }
        = \, \tikz[  x=1mm,
  y=1mm,
  baseline=-0.5ex] {
          \draw[->] (0,-6) -- (0,6) node [right] {$k$};
        }
      }.\label{diag-tag cancellation}
\end{align}
$\GS_{\zeta}$ applied to the LHS is just $\lambda_k\;\nu_{n-k}$ times $\GS_{\zeta}$ applied to the LHS. So this relation imposes the constraints $\lambda_{k} \;\nu_{n-k}=1$, which along with \cref{cond-tag switching} gives us that $\lambda_{k}=\nu_{k}^{-1}$, which is precisely \cref{cond-tag cancellation}.
It is easy to check that the mirror images and arrow reversals also yield the same conditions \cref{cond-tag cancellation}.

\subsection{I=H relation}\label{subsec- I=H}
For valid region labels, we have the I=H relations in $\cwebs$ given below (up to mirror images and arrow reversals):
\begin{gather}{
        \phantom{ \tikz[x=1mm, y=1mm, baseline=-0.5ex] {\draw (0,-10) -- (0,10)}} %for vertical spacing
        \tikz[x=1mm, y=1mm, baseline=-0.5ex] {
          \draw[->-] (-6,-6) node [left] {$l$}
            .. controls (-5,-5) and (-4,-3) .. (0,-3);
          \draw[->-] (6,-6) node [right] {$m$}
            .. controls (5,-5) and (4,-3) .. (0,-3);
          \draw[->-] (0,-3) -- (0,3);
          \draw[->-] (-6,6) node [left] {$k$}
            .. controls (-5,5) and (-5,3) .. (0,3);
          \draw[->] (0,3) .. controls (4,3) and (5,5) .. (6,6);
        }
        = \tikz[x=1mm, y=1mm, baseline=-0.5ex] {
          \draw[->-] (-6,-6) node [left] {$l$}
            .. controls (-5,-5) and (-3,-4) .. (-3,0);
          \draw[->-] (-6,6) node [left] {$k$}
            .. controls (-5,5) and (-3,4) .. (-3,0);
          \draw[->-] (-3,0) -- (3,0);
          \draw[->-] (6,-6) node [right] {$m$}
            .. controls (5,-5) and (3,-4) .. (3,0);
          \draw[->] (3,0) .. controls (3,4) and (5,5) .. (6,6);
        }
      }.
\end{gather}
Adding a clockwise cap on the top left, this relation is equivalent to
\begin{align}\label{diag-I=H relation}
    \isvg{0.15}{I=H_LHS} =\isvg{0.15}{I=H_RHS}.
\end{align}

So we desire the following equality in $\dmSBSBim$, for $k+l+m<n$:
\begin{align}\label{diag-image of I=H}
    \lambda_{k,l+m}\;\lambda_{l,m}\isvg{0.16}{image_of_I=H_LHS} =\lambda_{k,l}\;\lambda_{k+l,m}\isvg{0.16}{image_of_I=H_RHS}
\end{align}\vspace{5pt}\\
and the following one for $k+l+m=n$:
\begin{align}\label{diag-image of I=H k+l+m=n}
   \lambda_k \lambda_{l,m} \isvg{0.16}{image_of_I=H_LHS_k+l+m=n}= \lambda_{k,l} \lambda_{m} \isvg{0.16}{image_of_I=H_RHS_k+l+m=n}.
\end{align}

If the diagrams (without coefficients) appearing on either side of $\cref{diag-image of I=H}, \cref{diag-image of I=H k+l+m=n}$ were equal, then the desired relation is equivalent to the equality of the coefficients, which is precisely \cref{cond-I=H relation} (with the convention in \cref{defn convention for coefficients of trivalent with labels n} for $k+l+m=n$ case).
For \cref{diag-image of I=H k+l+m=n}, the equality of the diagrams follows from isotopy. For \cref{diag-image of I=H}, we prove it in the following lemma.

%For the edge case $n=k+l+m$, we have that the diagrams (without the coefficients) in the LHS and RHS of \cref{diag-image of I=H k+l+m=n} are the same up to isotopy. Hence \cref{diag-image of I=H k+l+m=n} holds if and only if
%$\lambda_{k}\lambda_{l,m}=\lambda_{k,l}\lambda_m$ when $k+l+m=n$. 
%Note that we may take $\lambda_{p}$ to be $\lambda_{p, n-p}$ so that this condition is already captured by \cref{cond-I=H relation}.

\begin{lemma}\label{lemma associativity ssbim}
\[ \isvg{0.16}{image_of_I=H_LHS}=\isvg{0.16}{image_of_I=H_RHS}\] in $\pdmSBSBim$, hence in $\dmSBSBim$.    
\end{lemma}
\begin{proof}
By color symmetry, we may assume $a=0$. 
With this choice, red represents $0$, blue $m$, green $l+m$ and yellow $k+l+m$. Note that when red and green are absent, blue and yellow are in different connected components and vice versa.

Using the distant Reidemeister II relation, we may cross the blue and yellow strands within the red-green bigon, as shown below.
\begin{equation}
    \isvg{0.15}{associativity_ssbim_LHS} = \isvg{0.15}{associativity_ssbim_1}
\end{equation}.

Then we do two oriented RIII relations to bring the yellow-blue crossings out of the red-green bigon.
% We first move the green strand past the yellow-blue crossing, and then move the red strand past the yellow-blue crossing, as shown below. 
\begin{equation}
 \isvg{0.15}{associativity_ssbim_1} = \isvg{0.15}{associativity_ssbim_2}=   \isvg{0.15}{associativity_ssbim_3}
\end{equation}
Now two more oriented RIII will bring the green-red crossings into the blue-yellow bigon.
% We first move the blue strand past the green-red crossing, and then move the yellow strand past the green-red crossing, as shown below.
\begin{equation}
    \isvg{0.15}{associativity_ssbim_3}=\isvg{0.15}{associativity_ssbim_4}=\isvg{0.15}{associativity_ssbim_5}
\end{equation}

Finally, we do a distant RII to uncross the green and red strands within the blue-yellow bigon.
\begin{equation}
  \isvg{0.15}{associativity_ssbim_5}=\isvg{0.15}{associativity_ssbim_6}.  
\end{equation}
The last term above is the RHS in \cref{lemma associativity ssbim} and we are done.
\end{proof}


Similarly, the relation obtained by taking the arrow reversals of \cref{diag-I=H relation} implies $\nu_{k, l+m}\;\nu_{l,m}=\nu_{k+l,m}\;\nu_{k,l}$.


\subsection{Relations and symmetries}

Above we checked specific versions of tag cancellation and $I=H$, and discussed some (but not all) of their mirror images and arrow reversals. This was enough to force the constraints on the scalars $\lambda$ and $\nu$ which were needed for \cref{lemma symmetry of GS-zeta}. Thanks to this lemma we do not need to check rotations, color changes via $\sigma$, or versions of relations obtained via $\Upsilon$.

For the remaining relations in $\cwebs$, namely \cref{diag-digon bursting} and \cref{diag-square flop} below, we only check one version of the relation. 

However, we must also verify the mirror images (resp. arrow reversals) of all the relations in $\cwebs$, and this does not correspond to a symmetry in $\DiagSBS$. Thankfully, the remaining mirror images and arrow reversals can be 
obtained alternatively from ones we have checked using only rotation and $\Upsilon$.

\subsection{Bursting a digon}\label{subsec-bigon bursting}
For valid region labels, we have the following ``bursting a digon" relation in $\cwebs$ (up to mirror images and arrow reversals).
\begin{align}\label{diag-digon bursting}{
        \tikz[x=1mm, y=1mm, baseline=-0.5ex] {
          \draw[->-] (0,-8) -- (0,-4);
          \draw[->-] (0,-4) arc (270:90:4);
          \draw[->-] (0,-4) arc (-90:90:4) node[pos=0.5,right]{$1$};
          \draw[->] (0,4) -- (0,8) node [right] {$k$};
        }
        = [k] \: \tikz[x=1mm, y=1mm, baseline=-0.5ex] {
          \draw[->] (0,-8) -- (0,8) node [right] {$k$};
        }
      }
      .\end{align}

For $k \neq 1$ or $n$, the LHS of the above diagram with a region label of $a$ on the right is sent to the following diagram via $\GS_{\zeta}$: 
\begin{align}\label{diag-image of bigon}
    \lambda_{k-1,1}\;\nu_{k-1,1}\isvg{0.15}{image_of_bigon}\;.
\end{align}

By color symmetry, we may assume that $a=0$. 
Using \cref{R2unoriented2}, \cref{R2oriented} and \cref{ccwcirc}, we can simplify the diagram in \cref{diag-image of bigon} (without the coefficients) to 
\begin{align}\label{eq-bigon burst simplification}
    \isvg{0.15}{ssbim_bigon_simplify_1}=\isvg{0.15}{ssbim_bigon_simplify_2} = \isvg{0.15}{ssbim_bigon_simplify_3}.
\end{align}
By \cref{corollary- bigon bursting}, we know that $\dd_{\hh{0,k}}^{\hh{0,1,k}}\left(\mu_{\hh{1}}^{\hh{0,1},\hh{1,k}}\right)=(-1)^{k-1}\zeta^{-k(k-1)/2}q^{-(k-1)}[k]_q$.
The above discussion then shows that $\GS_{\zeta}$ sends the LHS of \cref{diag-digon bursting} to 
\[(-1)^{k-1}\zeta^{-k(k-1)/2}q^{-(k-1)}[k]_q\; \lambda_{k-1,1} \; \nu_{k-1,1} \isvg{0.15}{image_of_tag_c}.\]
Comparing this with the image of the RHS of \cref{diag-digon bursting}, we have the conditions $$ \lambda_{k-1,1} \; \nu_{k-1,1}=(-1)^{k-1}\zeta^{k(k-1)/2}q^{k-1},$$ for $k\neq 1,n$ which is precisely \cref{cond-bigon killing}.

Now we consider the edge cases. 
For $k=1$, there is nothing to prove; for $k=n$, we use tag cancellation \cref{diag-tag cancellation}, and see that the resulting equality is true in $\dmSBSBim$ as follows.\\
Observe that $\GS_{\zeta}$ sends the the LHS to
\begin{align}\label{diag- bigon killing k=n}
  (-1)^{n-1}  \isvg{0.1}{image_of_bigon_k=n}.
\end{align}
By color symmetry, we may assume $a=0$. 
Using relations \cref{cccirc}, \cref{ccwcirc} we can simplify \cref{diag- bigon killing k=n} to $(-1)^{n-1} \dd_{\hh{0}}^{\hh{0,1}}\left(\mu^{\hh{0,1}}_{\hh{1}} \right)$.
From \cref{lemma demazure adjacent}, we have that $\dd_{\hh{0}}^{\hh{0,1}}\left(\mu^{\hh{0,1}}_{\hh{1}} \right)=(-1)^{n-1}[n]_q$, so that \cref{diag- bigon killing k=n} further reduces to $[n]_q$ as desired.

\subsection{Bursting a Square}\label{subsec-square flop}
For valid region labels, we had the following``bursting a square" relation in $\cwebs$:
\begin{align}\label{diag-square flop}
    {
      \phantom{ \tikz[x=1mm, y=1mm, baseline=-0.5ex] {\draw (0,-10) -- (0,10)}} %for vertical spacing
      \tikz[x=1mm, y=1mm, baseline=-0.5ex] {
        \draw[->-] (45:4) arc (45:135:4) node [pos=0.5, above]{$1$};
        \draw[-<-] (135:4) arc (135:225:4);
        \draw[->-] (225:4) arc (225:315:4) node [pos=0.5, below]{$1$};
        \draw[->-] (315:4) arc (315:405:4);
        \draw[->-] (7,-7) node [right] {$1$} -- (315:4);
        \draw[->] (45:4) -- (7,7) node [right] {$1$};
        \draw[->] (135:4) -- (-7,7) node [left] {$k$};
        \draw[->-] (-7,-7) node [left] {$k$} -- (225:4);
      }
      =
      \tikz[x=1mm, y=1mm, baseline=-0.5ex] {
        \draw[->-] (-7,-7) node [left] {$k$}
          .. controls (-6,-6) and (-3,-5) .. (0,-4);
        \draw[->-] (7,-7) node [right] {$1$}
          .. controls (6,-6) and (3,-5) .. (0,-4);
        \draw[->-] (0,-4) -- (0,4);
        \draw[->] (0,4) .. controls (-3,5) and (-6,6)
          .. (-7,7) node [left] {$k$};
        \draw[->] (0,4) .. controls (3,5) and (6,6)
          .. (7,7) node [right] {$1$};
      }
      + [k - 1] \tikz[x=1mm, y=1mm, baseline=-0.5ex] {
        \draw[->] (-7,-7) .. controls (-4,-4) and (-4,4)
          .. (-7,7) node [left] {$k$};
        \draw[->] (7,-7) .. controls (4,-4) and (4,4)
          .. (7,7) node [right] {$1$};
      }}.
\end{align}

The edge cases $k=0,1$ are trivial, and we treat the edge cases $k \in \{n-1,n\}$ afterwards, so assume $1 < k < n-1$. Applying $\GS_{\zeta}$, we then desire the following equality in $\dmSBSBim$: 
\begin{align}\label{diag-ssbim square flop}
   \lambda_{k-1,1} \nu_{k-1,1} \lambda_{1,1} \nu_{1,1} \isvg{0.15}{image_of_square_flop_LHS} = \lambda_{k,1}\nu_{k,1}\isvg{0.15}{image_of_square_flop_RHS_1} + [k-1]_q\isvg{0.15}{image_of_square_flop_RHS_2}.
\end{align}
The source and target of these morphisms is the object associated to the singlestep expression
\begin{equation} \label{flopsource} I_{\bullet} := [\hh{k+1} \supset \hh{1,k+1} \subset \hh{1} \supset \hh{0,1} \subset \hh{0}]. \end{equation}

From \cref{cond-bigon killing}, this simplifies to:
\begin{align}\label{diag-ssbim square flop simplified}
   (-1)^k\zeta^{\frac{k(k-1)}{2}+1}q^{k} \isvg{0.14}{image_of_square_flop_LHS} = (-1)^k\zeta^{\frac{k(k+1)}{2}}q^{k} \isvg{0.14}{image_of_square_flop_RHS_1} + [k-1]_q\isvg{0.14}{image_of_square_flop_RHS_2}.
\end{align}
We may check this in the algebraic category $\mSBSBim$, as an equality in the (degree zero) endomorphism ring of 
\[ \BS(I_{\bullet}) \cong \Bimod{\hh{k+1}}{}{R^{\hh{1,k+1}}}\otimes_{\hh{1}}\Bimod{}{\hh{0}}{R^{\hh{1,0}}} (\ell). \]
The grading shift $\ell$ is irrelevant when considering endomorphisms, and we omit it below.

Explicitly, the LHS of \cref{diag-ssbim square flop} corresponds to the $(R^{\hh{k+1}}, R^{\hh{0}})$-bimodule morphism given by
\begin{align}\label{eq 1 sec square flop}
    f\otimes g \mapsto (-1)^k\zeta^{\frac{k(k-1)}{2}+1}q^{k}\;\ddm{1,2, k+1}{1,k+1}\left(\ddm{1,2,k+1}{2,k+1}(\Del{1,2}{2}{1}f)\Del{1,2}{1}{1}\right) \otimes \ddm{0,1,2}{0,1}\left(\ddm{0,1,2}{0,2}(\Del{1,2}{2}{2}g)\Del{1,2}{1}{2} \right).
\end{align}
The first term on the RHS corresponds to
\begin{align}\label{eq 1.1 sec square flop}
    f\otimes g \mapsto (-1)^k\zeta^{\frac{k(k+1)}{2}}q^{k}\;\ddm{0,1,k+1}{1,k+1}\left( \ddm{0,1,k+1}{0, k+1}(fg)\Del{0,1}{1}{1}\right) \otimes \Del{0,1}{1}{2},
\end{align}
and the second term to
\begin{equation}\label{eq 1.2 sec square flop}
f \ot g \mapsto [k-1]_q(f\otimes g). \end{equation}

Checking that the map in  \cref{eq 1 sec square flop} agrees with the sum of the maps in \cref{eq 1.1 sec square flop} and \cref{eq 1.2 sec square flop} may be done after forgetting the bimodule structures and just working within the category of $A$-modules. Since $\BS(I_\bullet)$ is a free module over $A$ (since it free as a right module over $R^{\hh{0}}$ which is free over $A$), it suffices to check the desired equality after base extension from $A$ to its fraction field. 

We now proceed with the check after base extension to the fraction field of $A$. Evaluating the first term on the RHS at $1 \ot 1$ yields zero for degree reasons, and the equality of the LHS with the second term on the RHS is \cref{lemma square flop 1 tensor 1}. Evaluating all three terms at $1 \ot x_1^k$ is done in \cref{lemma square flop 1 tensor x1k}, confirming \cref{diag-ssbim square flop simplified} as applied to this element. Note that all three terms are nonzero when applied to $1 \ot x_1^k$. These two calculations also suffice to show that the two terms on the RHS are linearly independent as bimodule endomorphisms.

In \cref{lem:dim2} we prove that $\dim \End^0\left(\BS(I_{\bullet})\right)=2$. (Our $k$ here agrees with $k-1$ in \cref{lem:dim2}.)
Thus the two terms on the RHS of \cref{diag-ssbim square flop} are a basis for $\End^0$. The LHS of \cref{diag-ssbim square flop} must be a unique linear combination of the terms in the RHS. Checking the equality on $1 \ot 1$ and $1 \ot x_1^k$ gave two independent constraints for two unknowns, verifying \cref{diag-ssbim square flop}.

% \begin{lemma}\label{lemma dim of endomorphism ring is 2}
% For any $1\leq k \leq n-1$, let $I_{\bullet}$ be as in \eqref{flopsource}. Then  
% $$\dim \End^0\left(\BS(I_{\bullet}) \right)=2.$$
% \end{lemma}

% \begin{proof}
% \BE{I will rewrite other things so this works.}
% By the Soergel categorification theorem \BE{cref}, the character of $\BS(I_{\bullet})$ is $\bsh(I_{\bullet})$. By \BE{cref to lemma I need to write}, we have \begin{equation} \label{decatflop} \bsh(I_{\bullet}) = \KL_{\psi_0(\varpi_{k+1})} + \KL_{\psi_0(\varpi_1 + \varpi_k)}.\end{equation} By the asymptotic orthonormality of the KL basis \BE{cref}, this means that
% \begin{equation} (\bsh(I_{\bullet}), \bsh(I_{\bullet})) = 2 + v \Z[v], \end{equation}
% so by the Williamson-Soergel Hom formula \BE{cref}, we have $\dim \End^0(\BS(I_{\bullet})) = 2$.
% \end{proof}

% \begin{remark} On the other side of the Satake equivalence, following \BE{cref to unwritten remark}, this corresponds to $\dim \End(L_{\omega_k} \otimes L_{\omega_1}) = \dim \End(L_{\varpi_{k+1}} \oplus L_{\varpi_1 + \varpi_k}) = 2$. \end{remark}

% \BE{remove the below, but write some calculations earlier.}\PT{Actually, I did the full computation and we can indeed decompose the character of $M$ below as the Kazhdan-Lusztig basis for the coset $W_{\hh{k+1}}s_0s_n\ldots s_{k+1}W_{\hh{0}}$ and the coset of $1$, unless I made some mistake somewhere. Too late to double check and write it up here, but will add later today! The hom formula then immediately implies the result (along with asymptotic orthonormality of the KL basis).}
%     Let $M:=\Bimod{\hh{k+1}}{}{R^{\hh{1,k+1}}}\otimes_{\hh{1}}\Bimod{}{\hh{0}}{R^{\hh{1,0}}}.$
%     The bimodule $\Bimod{\hh{k+1}}{\hh{1}}{R^{\hh{1,k+1}}}$ coincides with the standard bimodule $\stdbimod{\hh{k+1}}{1}{\hh{1}}$ as defined in \cite[Definition 4.2.1]{WillSingular}.
%     Taking characters from \cref{thm-categorification}, we have that $$\sch{\hh{k+1}}{\hh{1}}([\stdbimod{\hh{k+1}}{1}{\hh{1}}])=v^{\ell(\hh{1,k+1})-\ell(\hh{1})}\;\;\HAB{\hh{k+1}}{\hh{1}}{1}$$ (see for example \cite[Lemma 2.45]{Abe}), which also equals  $v^{\ell(\hh{1,k+1})-\ell(\hh{1})}\;\;\KLB{\hh{k+1}}{\hh{1}}{1}$.
%     Similarly, we have 
%     $$\sch{\hh{1}}{\hh{0}}([\stdbimod{\hh{1}}{1}{\hh{0}}])=v^{\ell(\hh{0,1})-\ell(\hh{0})} \;\;\HAB{\hh{1}}{\hh{1}}{0}=v^{\ell(\hh{0,1})-\ell(\hh{0})}\;\;\KLB{\hh{1}}{\hh{0}}{1}.$$
%     So we have by \cref{thm-categorification} that 
%     $$\sch{\hh{k+1}}{\hh{0}} (M)=v^{\ell(\hh{0,1})+\ell(\hh{1,k+1})-\ell(\hh{0})-\ell(\hh{1})} \;\;\HAB{\hh{k+1}}{\hh{1}}{1}\ast_{\hh{1}}\HAB{\hh{1}}{\hh{0}}{1}.$$
%     The above term decomposes as a sum of two Kazhdan-Lusztig basis elements precisely as in \cref{example hecke conv square flop dimension computation}.
%     The Williamson-Soergel Hom formula then shows that 
%     $$\dim \End^\bullet (M)= 2$$ (from say, the asymptotic orthonormality of KL basis elements.)\PT{Add more details if needed!}
% %\end{proof}

Now we analyze the edge cases.

\subsubsection{When $k=n-1$}
In this case we desire the equality
\begin{align}\label{diag-square flop for k=n-1}
    \blkcoeff{\lambda_{n-2,1} \nu_{n-2,1} \lambda_{1,1} \nu_{1,1}} \isvg{0.15}{image_of_square_flop_LHS_for_n-1}=\blkcoeff{(-1)^{n-1}\lambda_{1}\nu_{n-1}}\isvg{0.15}{image_of_square_flop_RHS_1_for_n-1}+ \blkcoeff{[n-2]_q}\isvg{0.15}{image_of_square_flop_RHS_2_for_n-1}.
\end{align}
From \cref{cond-bigon killing}, \cref{cond-tag switching} and \cref{cond-tag cancellation}, we have that $\lambda_{n-2,1} \nu_{n-2,1} \lambda_{1,1} \nu_{1,1} =(-1)^k\zeta^{\frac{k(k+1)}{2}}q^{k}$ and $\lambda_1 \nu_{n-1}=1$.
As before, it suffices to check this equality in the algebraic category $\mSBSBim$, in the ring $\End\left(\Bimod{\hh{0}}{}{R^{\hh{0,1}}}\otimes_{\hh{1}}\Bimod{}{\hh{0}}{R^{\hh{1,0}}}\right).$
The LHS corresponds to the bimodule homomorphism in \cref{eq 1 sec square flop}, given by
\begin{align}
    f\otimes g \mapsto (-1)^{n-1}\zeta^{\frac{(n-1)(n-2)}{2}+1}q^{n-1}\;\ddm{0,1,2}{0,1}\left(\ddm{0,1,2}{0,2}(\Del{1,2}{2}{1}f)\Del{1,2}{1}{1}\right) \otimes \ddm{0,1,2}{0,1}\left(\ddm{0,1,2}{0,2}(\Del{1,2}{2}{2}g)\Del{1,2}{1}{2} \right)
\end{align}
while the RHS corresponds to
\begin{align}
    f\otimes g\mapsto (-1)^{n-1}\ddm{0,1}{0}(fg)\Del{0,1}{1}{1} \otimes \Del{0,1}{1}{2} \;+\: [n-2]_q(f\otimes g).
\end{align}

The proof is almost the same as the previous case, and we need only check that the desired equality holds after base extension to the fraction field of $A$. Evaluation of the LHS on $1 \ot 1$ is still done using \cref{lemma square flop 1 tensor 1}. Evaluation of all terms at $1 \ot x_1^{n-1}$ is done in \cref{lemma square flop 1 tensor x1k for k=n-1}. This checks \cref{diag-square flop for k=n-1} in two linearly-independent cases.
As before, we showed in \cref{lem:dim2} that $\dim \End\left(\Bimod{\hh{0}}{}{R^{\hh{0,1}}}\otimes_{\hh{1}}\Bimod{}{\hh{0}}{R^{\hh{1,0}}}\right)=2$, and that the terms on the right are linearly independent, so they should form a basis. Thus equality holds in \cref{diag-square flop for k=n-1}. 

\subsubsection{When $k=n$:}
In this case, the first term on the RHS of $\cref{diag-square flop}$ is $0$ and the term on the LHS may be simplified using tag cancellation \cref{diag-tag cancellation}, so that the relation becomes (with valid region labels) 
\begin{align}\label{diag-square flop k=n}
   {
        \tikz[x=1mm, y=1mm, baseline=-0.5ex] {
          \draw[->-] (0,-8) -- (0,-4);
          \draw[-<-] (0,-4) arc (270:90:4) node[pos=0.5, left]{$1$};
          \draw[->-] (0,-4) arc (-90:90:4) node[pos=0.5,right]{$2$};
          \draw[->] (0,4) -- (0,8) node [right] {$1$};
        }
        = [n-1] \: \tikz[x=1mm, y=1mm, baseline=-0.5ex] {
          \draw[->] (0,-8) -- (0,8) node [right] {$1$};
        }
      }.
\end{align}
The 2-functor $\GS_{\zeta}$ sends the LHS to 
\begin{align}\label{diag-image of square flop k=n}
    (-1)^{n-1}\lambda_{1,1}\nu_{1,1} \isvg{0.1}{image_of_square_flop_k=n} \;\;.
\end{align}
By color symmetry, we may assume that $a=0$. 
Using \cref{R2unoriented2}, \cref{R2oriented} and \cref{ccwcirc} as in \cref{eq-bigon burst simplification},
the diagram above (without the coefficients) simplifies to
\begin{align}\label{diag-square flop k=n simplify}
    \isvg{0.15}{ssbim_square_flop_n=k_simplify_1}=\isvg{0.15}{ssbim_square_flop_n=k_simplify_2}= \isvg{0.15}{ssbim_square_flop_n=k_simplify_3}.
\end{align}

We may simplify the polynomial in the box as follows.
\begin{align}\label{eq square flop k=n polynomial}
\notag \ddm{0,1,2}{0,1}\left(\mu^{\hh{0,2},\hh{1,2}}_{\hh{2}} \right) &= \zeta^{-\frac{(n-2)(n-3)}{2}}\dd_{n-1}\ldots \dd_3 \dd_2\left((x_3-\zeta^{n-1}x_2)(x_4- \zeta^{n-2}x_2)\ldots (x_n-\zeta^2x_2)\right)\\
\notag
    &\begin{array}{l}=\zeta^{-\frac{(n-2)(n-3)}{2}}(-1)^{n-2} \zeta^{(n-1)+      (n-2)+\ldots +2}\vspace{5 pt}\\
    \quad \quad\dd_{n-1}\ldots \dd_3\dd_2 \left((x_2-\zeta^{-2}x_n)(x_2-\zeta^{-3}x_{n-1})\ldots(x_2-\zeta^{-(n-1)}x_3)  \right)\end{array}\\
    &\begin{array}{l}=\zeta^{-\frac{(n-2)(n-3)}{2}}\zeta^{\frac{n(n-1)}{2}-1}(-1)^{n-2}\vspace{5 pt}\\
    \quad \quad\dd_{n-1}\ldots \dd_3\dd_2 \left((x_2-\zeta^{-2}x_n)(x_2-\zeta^{-3}x_{n-1})\ldots(x_2-\zeta^{-(n-1)}x_3)  \right).\end{array}
\end{align}
\cref{lemma demazure for bigon} and color symmetry tells us that 
$$\dd_{n-1}\ldots \dd_3\dd_2 \left((x_2-\zeta^{-2}x_n)(x_2-\zeta^{-3}x_{n-1})\ldots(x_2-\zeta^{-(n-1)}x_3)  \right)=q^{n-2}\zeta^{\frac{(n-2)(n-3)}{2}}[n-1]_q$$
so that \cref{eq square flop k=n polynomial} further simplifies to
\begin{align*}
   \ddm{0,1,2}{0,1}\left(\mu^{\hh{0,2},\hh{1,2}}_{\hh{2}} \right)  
   &= \zeta^{-\frac{(n-2)(n-3)}{2}}\zeta^{\frac{n(n-1)}{2}-1}(-1)^{n-2}q^{n-2}\zeta^{\frac{(n-2)(n-3)}{2}}[n-1]_q\\
   &=(-1)^{n-2}q^{-(n-1)}\zeta^{-1}q^{n-2}[n-1]_q =(-1)^{n-2}\zeta^{-1}q^{-1}[n-1]_q.
\end{align*}
Hence \cref{diag-image of square flop k=n} simplifies to
\[(-1)^{n-1} \lambda_{1,1}\nu_{1,1}\;(-1)^{n-2}\zeta^{-1}q^{-1}[n-1]_q=-\lambda_{1,1}\nu_{1,1}\zeta^{-1}q^{-1}[n-1]_q.\]
Comparing with the image of the RHS of \cref{diag-square flop k=n}, we have the conditions $\lambda_{1,1}\nu_{1,1}=-\zeta q.$
This is already captured in \cref{cond-bigon killing} for $k=2$.
